\section{The PA Role}
\textit{Reference $|$ Course Text: Chapters 1, 11, 13, 17}

At other institutions, the PA role may be called ``faculty mentor,'' ``project advisor,'' ``capstone instructor,'' or ``section leader.'' Regardless of title, the function is the same: you are the primary point of contact between the course and the team, and your judgment is the most influential factor in each student's experience.

\subsection{Core Responsibilities}

\begin{enumerate}[nosep,leftmargin=1.5em]
\item \textbf{Mentoring.} Guide teams through the design process without doing the design for them. Ask questions that expose gaps in reasoning. Suggest approaches without prescribing solutions.

\item \textbf{Grading and Assessment.} Provide PA input on team and individual grades. Evaluate deliverables for engineering quality, not just completeness. Complete Individual Performance Evaluations on schedule.

\item \textbf{Gate-Keeping.} Verify that teams meet minimum readiness criteria before advancing through design reviews (PDR, CDR, FDR). A team that is not ready for CDR should not present at CDR.

\item \textbf{Client Liaison.} Maintain awareness of the client relationship. Intervene if communication breaks down. Attend client meetings when invited or when the relationship needs support.

\item \textbf{Safety Oversight.} Review and approve Hazard Assessments before fabrication begins. Escalate safety concerns to the CLT Director immediately.
\end{enumerate}

\subsection{Relationship with CF, CLT Director, and Other PAs}

\begin{itemize}[nosep,leftmargin=1.5em]
\item The \textbf{Course Faculty (CF)} sets course policy, owns the syllabus, and makes final grade determinations. PAs provide input; the CF decides.
\item The \textbf{CLT Director} manages program logistics, client relationships at the organizational level, and facilities (InnoHub, labs). Escalate facility, safety, or client-relationship issues to the CLT Director.
\item \textbf{Other PAs} are your peers. Calibrate grading expectations by comparing notes at PA meetings. If you are unsure whether a deliverable meets the bar, ask a fellow PA for a second opinion.
\end{itemize}

\subsection{Time Commitment}

Expect to invest approximately \textbf{3--5 hours per week per team}, distributed roughly as follows:

{\centering\small
\begin{tabularx}{\textwidth}{Xr}
\toprule
\textbf{Activity} & \textbf{Approx.\ Hours/Week} \\
\midrule
Sprint review meetings (preparation + attendance) & 1.0--1.5 \\
Deliverable review and feedback & 1.0--1.5 \\
Email, Slack, and ad-hoc student questions & 0.5--1.0 \\
PA coordination meetings and calibration & 0.25--0.5 \\
Design review preparation (PDR/CDR/FDR weeks) & +2.0--3.0 \\
\bottomrule
\end{tabularx}
\par}

\vspace{4pt}
Time investment increases during design review sprints and decreases during steady-state execution sprints (Sprints 9--11).

\subsection{Authority and Boundaries}

\begin{faqbox}[What does the PA decide vs.\ what does the CF decide?]
\textbf{PA decides:}
\begin{itemize}[nosep,leftmargin=1.5em]
\item Whether a team is ready for a design review (gate-keeping)
\item Day-to-day mentoring approach and meeting cadence
\item Sprint review feedback and coaching priorities
\item Input on individual and team grade recommendations
\end{itemize}

\textbf{CF decides:}
\begin{itemize}[nosep,leftmargin=1.5em]
\item Final grades and grade disputes
\item Course policy exceptions (late work, extensions, accommodations)
\item Academic integrity cases
\item Team reassignments or project changes
\item Performance Improvement Plans (PIPs)
\end{itemize}
\end{faqbox}

\begin{warnbox}
Do not make promises about grades. You may say ``this is strong A-level work'' or ``this needs significant improvement to pass,'' but final grade determination is the CF's responsibility. Overpromising grades creates conflicts that are difficult to resolve.
\end{warnbox}

\begin{bestpractice}
Establish your communication norms early. Tell your teams how you prefer to receive questions (email, Slack, office hours), your typical response time, and what constitutes an emergency worth a same-day message. Setting expectations in Sprint~0 prevents friction later.
\end{bestpractice}


% ═══════════════════════════════════════
